\vspace{-3mm}
\section{Problem Formulation}
\label{sec:ProblemForm}
\vspace{-2mm}
Consider $N$ agents in the scene. Let $\mathcal{X}_i$ and $\mathcal{Y}_i$ be the observation and the perception supervision of the $i$th agent, respectively. The objective of collaborative perception is to achieve the maximized perception performance of all agents as a function of the total communication budge $B$ and communication round $K$; that is,
\vspace{-1mm}
\begin{equation*}
    {\xi}_{\Phi}(B, K) = \underset{\theta,\mathcal{P}}{\arg \max}~\sum_{i=1}^{N} 
    g \left(\Phi_{\theta} \left(\mathcal{X}_i,\{\mathcal{P}_{i\rightarrow j}^{(K)}\}_{j=1}^N \right), \mathcal{Y}_i  \right),~
    ~{\rm s.t.~} \sum_{k=1}^{K} \sum_{i=1}^{N}|\mathcal{P}_{i\rightarrow j}^{(k)}| \leq B,
\vspace{-1mm}
\end{equation*}
where $g(\cdot,\cdot)$ is the perception evaluation metric, $\Phi$ is the perception network with trainable parameter $\theta$, and $\mathcal{P}_{i\rightarrow j}^{(k)}$ is the message transmitted from the $i$th agent to the $j$th agent at the $k$th communication round. Note that i) when $B=K=0$, there is no collaboration and ${\xi}_{\Phi}(0, 0)$ reflects the single-agent perception performance; ii) through optimizing the communication strategy and the network parameter, collaborative perception should perform well consistently at any communication bandwidth or round; and iii) we consider multi-round communication, where each agent serves as both a supporter (offering message to help others) and a requester (requesting messages from others).



In this work, we consider the perception task of 3D object detection and present three contributions: i) we make communication more efficient by designing compact messages and sparse communication graphs; ii) we boost the perception performance by implementing more comprehensive message fusion; iii) we enable the overall system to adapt to varying communication conditions by dynamically adjusting where and who to communicate.
% {\HC we enable the overall system to adapt to varying communication conditions by strategically allocating communication budges based on the perceptually spatial criticalness.}
% \textcolor{red}{spatially perceptual significance}.


\vspace{-2mm}
\section{\texttt{Where2comm}: Spatial Confidence-Aware Collaborative Perception System}
\vspace{-2mm}

This section presents~\texttt{Where2comm}, a multi-round, multi-modality, multi-agent collaborative perception framework based on a spatial-confidence-aware communication strategy; see the overview in Fig.~\ref{fig:system}. \texttt{Where2comm} includes an observation encoder, a spatial confidence generator, the spatial confidence-aware communication module, the spatial confidence-aware message fusion module and a detection decoder. Among five modules, the proposed spatial confidence generator generates the spatial confidence map. Based on this spatial confidence map, the proposed spatial confidence-aware communication generates compact messages and sparse communication graphs to save communication bandwidth; and the proposed spatial confidence-aware message fusion module leverages informative spatial confidence priors to achieve better aggregation; also see an algorithmic summary in Algorithm 1 and the optimization-oriented design rationale in Section 7.3 in Appendix.


\begin{figure}[!t]
    \centering
    % \includegraphics[width=1.0\linewidth]{figures/system-0518.png}
    \includegraphics[width=1.0\linewidth]{figures/System/system-final.png}
    % \includegraphics[width=1\linewidth]{figures/system_v2.png}
    \caption{System overview. In \texttt{Where2comm}, spatial confidence generator enables the awareness of spatial heterogeneous of perceptual information, spatial confidence-aware communication enables efficient communication, and spatial confidence-aware message fusion boosts the performance.}
    \label{fig:system}
    \vspace{-5mm}
\end{figure}

\vspace{-2mm}
\subsection{Observation encoder}
\label{sec:Encoder}
\vspace{-2mm}
The observation encoder extracts feature maps from the sensor data.~\texttt{Where2comm} accepts single/multi-modality inputs, such as RGB images and 3D point clouds. This work adopts the feature representations in bird's eye view (BEV), where all agents project their individual perceptual information to the same global coordinate system, avoiding complex coordinate transformations and supporting better shared cross-agent collaboration. For the $i$th agent, given its input $\mathcal{X}_i$, the feature map is 
$ 
\mathcal{F}_i^{(0)} = \Phi_{\rm enc}(\mathcal{X}_i) \in \mathbb{R}^{H \times W \times D},$
where $\Phi_{\rm enc}(\cdot)$ is the encoder, the superscript $0$ reflects that the feature is obtained before communication and $H,W,D$ are its height, weight and channel. All agents share the same BEV coordinate system. For the image input, $\Phi_{\rm enc}(\cdot)$ is followed by a warping function that transforms the extracted feature from front-view to BEV. For 3D point cloud input, we discretize 3D points as a BEV map and $\Phi_{\rm enc}(\cdot)$ extracts features in BEV. The extracted feature map is output to the spatial confidence generator and the message fusion module.


\vspace{-2mm}
\subsection{Spatial confidence generator}
\label{sec:Generator}
\vspace{-2mm}
The spatial confidence generator generates a spatial confidence map from the feature map of each agent. The spatial confidence map reflects the perceptually critical level of various spatial areas. Intuitively, for object detection task, the areas that contain objects are more critical than background areas. During collaboration, areas with objects could help recover the miss-detected objects due to the limited view; and background areas could be omitted to save the precious bandwidth. So we represent the spatial confidence map with the detection confidence map, where the area with high perceptually critical level is the area that contains an object with a high confidence score.

To implement, we use a detection decoder structure to produce the detection confidence map. Given the feature map at the $k$th communication round, $\mathcal{F}_i^{(k)}$, the corresponding spatial confidence map is 
\vspace{-1mm}
\begin{equation}
\label{eq:generator}
    \mathbf{C}_i^{(k)}  =  \Phi_{\rm generator}(\mathcal{F}_i^{(k)})  \in [0,1]^{H \times W},
\end{equation}
%\vspace{-1mm}
where the generator $\Phi_{\rm generator}(\cdot)$ follows a detection decoder. Since we consider multi-round collaboration, \texttt{Where2comm} iteratively updates the feature map by aggregating information from other agents. Once $\mathcal{F}_i^{(k)}$ is obtained,~\eqref{eq:generator} is triggered to reflect the perceptually critical level at each spatial location. The proposed spatial confidence map answers a crucial question that was ignored by previous works: for each agent, information at which spatial area is worth sharing with others. By answering this, it provides a solid base for efficient communication and effective message fusion.

\vspace{-2mm}
\subsection{Spatial confidence-aware communication}
\label{sec:Comm}
\vspace{-2mm}
With the guidance of spatial confidence maps, the proposed communication module packs compact messages with spatially sparse feature maps and transmits messages through a sparsely-connected communication graph. Most existing collaboration perception systems~\cite{v2vnet,disconet,xu2022v2x} considers full feature maps in the messages and fully-connected communication graphs. To reduce the communication bandwidth without affecting perception, we leverage the spatial confidence map to select the most informative spatial areas in the feature map (where to communicate) and decide the most beneficial collaboration partners (who to communicate).

\textbf{Message packing.} Message packing determines what information should be included in the to-be-sent message. The proposed message includes: i) a request map that indicates at which spatial areas the agent needs to know more; and ii) a spatially sparse, yet perceptually critical feature map.

The request map of the $i$th agent is $\mathbf{R}_{i}^{(k)} = 1 - \mathbf{C}_{i}^{(k)} \in \mathbb{R}^{H\times W}$, negatively correlated with the spatial confidence map. The intuition is, for the locations with low confidence score, an agent is hard to tell if there is really no objects or it is just caused by the limited information (e.g. occlusion). Thus, the low confidence score indicates there could be missing information at that location. Requesting information at these locations from other agents could improve the current agent's detection accuracy.

The spatially sparse feature map are selected based on each agent's spatial confidence map and the received request maps from others. Specifically, a binary selection matrix is used to represent each location is selected or not, where $1$ denotes selected, and $0$ elsewhere. For the message sent from the $i$th agent to the $j$th agent at the $k$th communication round, the binary selection matrix is
\begin{equation}
\vspace{-1mm}
\label{eq:selector}
\mathbf{M}_{i\rightarrow j}^{(k)}=\left\{
\begin{array}{lll}
\Phi_{\rm select}(\mathbf{C}_{i}^{(k)})\in \{0, 1\}^{H\times W}, &\quad k=0;\\
\Phi_{\rm select}(\mathbf{C}_{i}^{(k)}\odot \mathbf{R}_j^{(k-1)} ), \in \{0, 1\}^{H\times W},  &\quad k>0;
\vspace{-1mm}
\end{array} \right.
\end{equation}
where $\odot$ is the element-wise multiplication, $\mathbf{R}_j^{(k-1)}$ is the request map from the $j$th agent received at the previous round, $\Phi_{\rm select}(\cdot)$ is the selection function which targets to select the most critical areas conditioned on the input matrix, which represents the critical level at the certain spatial location. We implement $\Phi_{\rm select}(\cdot)$ by selecting the locations where the largest elements at in the given input matrix conditioned on the bandwidth limit; optionally, a Gaussian filter could be applied to filter out the outliers and introduce some context. In the initial communication round, each agent selects the most critical areas from its own perspective as the request maps from other agents are not available yet; in the subsequent rounds, each agent also takes the partner's request into account, enabling more targeted communication. Then, the selected feature map is obtained as
$
    \mathcal{Z}_{i\rightarrow j}^{(k)}=\mathbf{M}_{i\rightarrow j}^{(k)}\odot \mathcal{F}_i^{(k)} \in \mathbb{R}^{H\times W \times D},
$
which provides spatially sparse, yet perceptually critical information. 

Overall, the message sent from the $i$th agent to the $j$th agent at the $k$th communication round is $\mathcal{P}_{i\rightarrow j}^{(k)}=(\mathbf{R}_{i}^{(k)}, \mathcal{Z}_{i\rightarrow j}^{(k)} )$. Note that i) $\mathbf{R}_{i}^{(k)}$ provides spatial priors to request complementary information for the $i$th agent's need in the next round; the feature map $ \mathcal{Z}_{i\rightarrow j}^{(k)}$ provides supportive information for the $i$th agent's need in the this round. They together enable mutually beneficial collaboration; ii) since  $\mathcal{Z}_{i\rightarrow j}^{(k)}$ is sparse, we only transmit non-zero features and corresponding indices, leading to low communication cost; and iii) the sparsity of $\mathcal{Z}_{i\rightarrow j}^{(k)}$ is determined by the binary selection matrix, which dynamically allocates the communication budget at various spatial areas based on their perceptual critical level, adapting to various communication conditions. 




\textbf{Communication graph construction.} Communication graph construction targets to identify when and who to communicate to avoid unnecessary communication that wastes the bandwidth. Most previous works~\cite{v2vnet,disconet,OPV2V} consider fully-connected communication graphs. When2com~\cite{when2com} proposes a handshake mechanism, which uses similar global features to match partners. This is hard to interpret because two agents, which have similar global features, do not necessarily need information from each other. Different from all previous works, we provide an explicit design rationale: the necessity of communication between the $i$th and the $j$th agents is simply measured by the overlap between the information that the $i$th agent has and the information that the $j$th agent needs. With the help of the spatial confidence map and the request map, we construct a more interpretable communication graph.

For the initial communication round, every agent in the system is not aware of other agents yet. To activate the collaboration, we construct a fully-connected communication graph. Every agent will broadcast its message to the rest of the system. For the subsequent communication rounds, we examine if the communication between agent $i$ and agent $j$ is necessary based on the maximum value of the binary selection matrix $\mathbf{M}_{i\rightarrow j}^{(k)}$, i.e. if there is at least one patch is activated, then we regard the connection is necessary. Formally, let $\mathbf{A}^{(k)}$ be the adjacency matrix of the communication graph at the $k$th communication round, whose $(i,j)$th element is \vspace{-1mm}
\begin{equation*}
 \mathbf{A}_{i,j}^{(k)}=\left\{
\begin{array}{lll}
1, &\quad k=0; \\
{\rm max}_{h \in \{0,1,..,H-1\},w \in \{0,1,...,W-1\}}~ \left( \mathbf{M}_{i\rightarrow j}^{(k)} \right)_{h,w} \in \{0, 1\}, &\quad k>0;
\end{array} \right.
\end{equation*}
where $h, w$ index the spatial area, reflecting message passing from the $i$th agent to the $j$th agent.  Given this sparse communication graph, agents can exchange messages with selected partners.

% \textbf{A selectivelydvantages summary} Compared to the existing methods, the proposed fine-grained spatial-aware communication i) improves the bandwidth usage efficiency by enabling agents to selectively pack the most critical feature patches; ii) ensures the communication effectiveness by allowing agents to easily convey its needs in the collaboration through a request map; iii) constructs communication graphs in an much more interpretable way.

\vspace{-2mm}
\subsection{Spatial confidence-aware message fusion}
\label{sec:Fusion}
\vspace{-2mm}
Spatial confidence-aware message fusion targets to augment the feature of each agent by aggregating the received messages from the other agents. To achieve this, we adopt a transformer architecture, which leverages multi-head attention to fuse the corresponding features from multiple agents at each individual spatial location. The key technical design is to include the spatial confidence maps of all the agents to promote cross-agent attention learning. The intuition is that, the spatial confidence map could explicitly reflect the perceptually critical level, providing a useful prior for attention learning.

Specifically, for the $i$th agent, after receiving the $j$th agent's message $\mathcal{P}_{j\rightarrow i}^{(k)}$, it could unpack to retrieve the feature map $\mathcal{Z}^{(k)}_{j\rightarrow i}$ and the spatial confidence map $\mathbf{C}^{(k)}_j = 1 - \mathbf{R}^{(k)}_j$. We also include the ego feature map in fusion and denote $\mathcal{Z}^{(k)}_{i\rightarrow i} = \mathcal{F}_{i}^{(k)}$ to make the formulation simple and consistent, where $\mathcal{Z}^{(k)}_{i\rightarrow i}$ might not be sparse. To fuse the features from the $j$th agent at the $k$th communication round, the cross-agent/ego attention weight for the $i$th agent is
\begin{equation}
    \label{eq:attention}
    \mathbf{W}_{j\rightarrow i}^{(k)} = {\rm MHA}_{\rm W}\left(\mathcal{F}_{i}^{(k)}, \mathcal{Z}_{j\rightarrow i}^{(k)}, \mathcal{Z}_{j\rightarrow i}^{(k)}\right) \odot  \mathbf{C}^{(k)}_j \in \mathbb{R}^{H\times W},
\end{equation}
where ${\rm MHA}_{\rm W}(\cdot)$ is a multi-head attention applied at each individual spatial location, which outputs the scaled dot-product attention weight. Note that i) the proposed spatial confidence maps contributes to the attention weight, as the features with higher perceptually critical level are more preferred in the feature aggregation; ii) the cross-agent attention weight models the collaboration strength with a $H\times W$ spatial resolution, leading to more flexible information fusion at various spatial regions. Then, the feature map of the $i$th agent after fusing the messages in the $k$th communication round is
\begin{align*}
\vspace{-4mm}
    \mathcal{F}_{i}^{(k+1)} &= {\rm FFN} \left( \sum_{j\in  \mathcal{N}_i\bigcup \{i\} } {\mathbf{W}}_{j\rightarrow i}^{(k)} \odot \mathcal{Z}_{j\rightarrow i}^{(k)}   \right) \in \mathbb{R}^{H\times W\times D},
\vspace{-4mm}
\end{align*}
where ${\rm FFN}(\cdot)$ is the feed-forward network and $\mathcal{N}_i$ is the neighbors of the $i$th agent defined in the communication graph $\mathbf{A}^{(k)}$. The fused feature $\mathcal{F}_{i}^{(k+1)}$ would serve as the $i$th agent's feature in the $(k+1)$th round. In the final round, we output $\mathcal{F}_{i}^{(k+1)}$ to the detection decoder to generate detections.

\textbf{Sensor positional encoding.} Sensor positional encoding represents the physical distance between each agent's sensor and its observation.  It adopts a standard positional encoding function conditioned on the sensing distance and feature dimension. The features are summed up with the positional encoding of each location before inputting to the transformer.

Compared to existing fusion modules that do not use attention mechanism~\cite{v2vnet} or only use agent-level attentions~\cite{when2com}, the per-location attention mechanism adopted by the proposed fusion emphasizes the location-specific feature interactions. It  makes the feature fusion more targeted. Compared to the methods that also use the per-location attention-based fusion module\cite{disconet,OPV2V,xu2022v2x}, the proposed fusion module leverages multi-head attention with two extra priors, including spatial confidence map and sensing distances. Both assist attention learning to prefer high quality and critical features.

\vspace{-2mm}
\subsection{Detection decoder}
\label{sec:Decoder}
\vspace{-2mm}
The detection decoder decodes features into objects, including class and regression output. 
Given the feature map at the $k$th communication round $\mathcal{F}_i^{(k)}$, the detection decoder $\Phi_{\rm dec}(\cdot)$ generate the detections of $i$th agent by
% The detections of $i$th agent are
$
     \widehat{\mathcal{O}}_i^{(k)}  =  \Phi_{\rm dec}(\mathcal{F}_i^{(k)}) \in \mathbb{R}^{H \times W \times 7},
$
% where $\mathcal{F}_i^{(k)}$ is the feature map at the $k$th communication round, $\Phi_{\rm dec}(\cdot)$ is the detection decoder, $\widehat{\mathcal{O}}_i^{(k)}$ is the detection map, 
where each location of $\widehat{\mathcal{O}}_i^{(k)}$ represents a rotated box with class $(c,x,y,h,w, \cos\alpha, \sin\alpha)$, denoting class confidence, position, size and angle. The objects are the final output of the proposed collaborative perception system. Note that $\widehat{\mathcal{O}}_i^{(0)}$ denotes the detections without collaboration.

\vspace{-2mm}
\subsection{Training details and loss functions}
\vspace{-2mm}
To train the overall system, we supervise two tasks: spatial confidence generation and object detection at each round. As mentioned before, the functionality of the spatial confidence generator is the same as the classification in the detection decoder. To promote parameter efficiency, our spatial confidence generator reuses the parameters of the detection decoder. For the multi-round settings, each round is supervised with one detection loss, the overall loss is
$
    L = \sum_{k=0}^K\sum_i^{N} L_{\rm det} \left(\widehat{\mathcal{O}}_i^{(k)},\mathcal{O}_i \right),
$
where $\mathcal{O}_i$ is the $i$th agent's ground-truth objects, $L_{\rm det}$ is the detection loss~~\cite{zhou2019objects}.


\textbf{Training strategy for multi-round setting.} To adapt to multi-round communication and dynamic bandwidth, we train the model under various communication settings with curriculum learning strategy~\cite{curriculum}. We first gradually increase the communication bandwidth and round; and then, randomly sample bandwidth and round to promote robustness. Through this training strategy, a single model can perform well at various communication conditions.
%  Through this training strategy, the model can adapt to learns how to allocate communication budges at various communication conditions.
% We allocate a small bandwidth (about 20\%) for the initial redundant broadcast communication round; then a relatively large (about 60\%) bandwidth for the next targeted and efficient communication round enabled by agents' request; afterwards gradually reduce the bandwidth considering communication degradation with the increasing rounds. 